\begin{definition}[Viscosity supersolution of $L(\nabla u,\nabla^2u)=f$, paper Def.~def-vis]
Let $n\ge0$, $\Omega\subseteq\mathbb R^n$, $L:\mathbb R^n\times\mathbb R^{n\times n}\to\mathbb R$ and $f,u:\mathbb R^n\to\mathbb R$. We say that $u$ is a \emph{viscosity supersolution of $L(\nabla u,\nabla^2u)=f$ in $\Omega$} if
\begin{enumerate}
\item[(i)] $u$ is lower semicontinuous on $\Omega$ (relative to $\Omega$: for every $x\in\Omega$ and $t<u(x)$ there is a neighbourhood $U$ of $x$ with $u>t$ on $U\cap\Omega$), and
\item[(ii)] whenever $x_0\in\Omega$ and $\psi:\mathbb R^n\to\mathbb R$ is $C^2$ on $\Omega$ (Lean \texttt{ContDiffOn ℝ 2 ψ Ω}) are such that $u-\psi$ attains a local minimum at $x_0$ (i.e.\ $u(x)-\psi(x)\ge u(x_0)-\psi(x_0)$ for all $x$ in some neighbourhood of $x_0$ in $\mathbb R^n$), we have
\[
f(x_0)\le L\bigl(\nabla\psi(x_0),\nabla^2\psi(x_0)\bigr),
\]
where $\nabla\psi(x_0)$ is the gradient and $\nabla^2\psi(x_0)$ the Hessian matrix of \noderef{EuclidHessian}.
\end{enumerate}
For an open $\Omega$ this is exactly the paper's Definition def-vis of a viscosity supersolution $u\in LSC(\Omega)$ of $L(\nabla u,\nabla^2u)=f$, i.e.\ of $\tilde L(x,\xi,X)=L(\xi,X)-f(x)=0$ ($\tilde L\ge0$ at test points), with test functions $\psi\in C^2(\Omega)$.
\end{definition}
